
\documentclass[11pt]{article}
\usepackage[margin=1in]{geometry}
\usepackage{amsmath,amssymb,booktabs,hyperref,graphicx,parskip,enumitem}
\usepackage{xcolor}
\hypersetup{colorlinks=true,linkcolor=blue!50!black,urlcolor=blue!50!black}
\setlist{nosep}

\title{Replication Report: Quantitative Relationship between Polarization Differences and the Zone-Averaged Shift Photocurrent\\[2pt]\large OSTI 1523841 \textbar{} Coverage 8/10, Agreement 10/10}
\author{Rick Stevens \and Ollie (AI Assistant)\\\small Argonne National Laboratory}
\date{2026-04-27}

\begin{document}\maketitle

\section{Source paper}
\begin{itemize}
\item \textbf{Title:} Quantitative Relationship between Polarization Differences and the Zone-Averaged Shift Photocurrent
\item \textbf{Authors:} Fregoso, Morimoto, Moore
\item \textbf{Year:} 2017
\item \textbf{Domain:} Condensed Matter / Topological Response / Shift Current Photovoltaics
\item \textbf{Identifier:} OSTI 1523841
\end{itemize}


\section{Replication objective}
The central claim of the paper concerns condensed matter / topological response / shift current photovoltaics. Our objective was to reproduce the headline numerical/qualitative results within the constraints of the AI-driven Replication Project (24--72\,h, commodity GPU/CPU compute, no author contact). A replication plan is archived at \texttt{1523841-Quantitative-relationship-between-polarization-differences-a/replication\_plan.pdf}.

\section{Methodology}
We followed the standard Replication Project workflow: ingest the source PDF, extract method and data pointers, draft a replication plan, attempt the smallest end-to-end pipeline that produces a comparable number, then scale up where compute and data permit. Tools used in this replication are catalogued in the meta-paper's computational tools inventory (see \texttt{COMPUTATIONAL\_TOOLS\_INVENTORY.md}).



\subsection*{Paper-directory README excerpt}
\small
\par\medskip\textbf{Quantitative relationship between polarization differences and the zone-averaged shift photocurrent}\par
\par$\bullet$\ **OSTI ID:** 1523841
\par$\bullet$\ **Rank:** \#22
\par$\bullet$\ **Replication Score:** 9/10
\par$\bullet$\ **Open-Source Tools:** Yes
\par$\bullet$\ **Code Repository:** No
\par$\bullet$\ **Tools:** None explicitly (tight-binding models, analytic equations; implementation possible in Python, Julia, etc.)

\par\medskip\textbf{\# Why This Paper}\par
Purely analytic/numerical, all equations and parameters specified, synthetic data, and quantitative results. No code repo but trivial to implement for AI.

\par\medskip\textbf{\# Replication Plan}\par
AI implements tight-binding models, computes Berry connections/shift vectors, and reproduces tables/curves using open-source tools.

\par\medskip\textbf{\# Status}\par
\par$\bullet$\ [ ] Paper reviewed
\par$\bullet$\ [ ] Code/tools identified
\par$\bullet$\ [ ] Code implemented/cloned
\par$\bullet$\ [ ] Results reproduced
\par$\bullet$\ [ ] Results validated against paper
\normalsize

The replication was driven by an OpenClaw agent loop (Claude Opus / GPT-5 class models) issuing shell, Python, and (where applicable) GPU jobs. Compute usage and wall-time for individual papers are summarised in the meta-paper's resource table; not separately recorded for this paper unless noted in the Tier-Lift artefact. Sub-agents were spawned for replication-plan generation and (for papers in batches 1--4) for tier-lift extension runs.

\section{What was reproduced}
\begin{itemize}
\item All previous items (Eq.9, Eq.C5, Eq.D4, Eq.D8-D9, winding discontinuity, Fig.1(b))
\item Multi-band Wilson-loop/Eq.17 verification on 3-band trimer
\item Multi-band test on 4-band coupled RM model
\item 1D BHZ-like spin-doubled model (W=0 trivial check)
\item Integer winding W\_cv for every band pair (gauge-independent)
\item Berry polarizations P\_n mod 1 invariance under parallel transport
\end{itemize}

\section{What was missing or incomplete}
\begin{itemize}
\item Material-specific DFT+Wannier (GeS, BaTiO3, monolayer WS2)
\item Frequency-integrated shift-conductivity connection σ\textasciicircum{}shift(ω)
\item Paper Figs 2-4 (higher-D extensions, material examples)
\item Comparison vs experimental shift-current measurements
\item Hartree-Fock/GW interacting-system extension
\end{itemize}
The dominant blocker categories for this paper, mapped onto the meta-paper's 9-category friction taxonomy (\S4), are listed in \S\ref{sec:friction} below.

\section{Quantitative agreement}
Per-quantity numerical comparisons are not separately tabulated in the structured evaluation record for this paper beyond the agreement rationale below; where the rationale references specific numbers, they are reproduced verbatim. A full numerical comparison table is recorded in the per-replication artefacts (see \S\ref{sec:repro}). For papers below 8/8, the agreement rationale identifies the specific quantities that diverge.

\paragraph{Agreement rationale.} Two-band identity continues to hold to machine precision on the Rice-Mele model. New multi-band tests confirm gauge-invariant polarizations and integer winding numbers (BHZ: W=0 trivially; trimer/RM-block: integer winding tracked through level reordering). Identity (P\_c-P\_v)+W\_cv is preserved as the universal structure across all tested band pairs.

\section{Score and friction-point tagging}\label{sec:friction}
\begin{itemize}
\item \textbf{Coverage:} 8/10 --- Original c=7 lifted to 8 by adding a multi-band tier-lift extension (code/multiband\_extension.py) that verifies the polarization-difference identity beyond the two-band Rice-Mele case. We tested the Wilson-loop/Eq.17 generalisation on (1) a 3-band trimer chain, (2) a 4-band coupled Rice-Mele block model, and (3) a 1D BHZ-like spin-doubled model. Berry polarizations are gauge-invariant mod 1, the dipole-phase winding W\_cv is integer for every band pair, and the residual P\_c-P\_v+W\_cv is well-defined. This addresses follow-on Q1 directly. Still missing: material-specific DFT+Wannier calculations (GeS, BaTiO3) and the frequency-integrated shift-conductivity connection (Q2/Q3).
\item \textbf{Agreement:} 10/10 --- Two-band identity continues to hold to machine precision on the Rice-Mele model. New multi-band tests confirm gauge-invariant polarizations and integer winding numbers (BHZ: W=0 trivially; trimer/RM-block: integer winding tracked through level reordering). Identity (P\_c-P\_v)+W\_cv is preserved as the universal structure across all tested band pairs.
\end{itemize}

\noindent\textbf{Friction points encountered (taxonomy tags):}
\begin{itemize}
\item \textbf{F7}: Missing surrogate calculators (xTB/DFT/closed solver)
\end{itemize}

\section{Follow-on questions}
\begin{itemize}
\item Q2/Q3/Q5 from prior entry remain open
\end{itemize}

\section{Reproducibility pointers}\label{sec:repro}
\begin{itemize}
\item \textbf{Paper directory:} \texttt{1523841-Quantitative-relationship-between-polarization-differences-a}
\item \textbf{Replication artefacts:}
\begin{itemize}
\item \texttt{/Users/stevens/Dropbox/REPLICATE-PROJECT/1523841-Quantitative-relationship-between-polarization-differences-a/replication}
\end{itemize}
\item \textbf{Replication-dir hint from JSONL:} \texttt{\textasciitilde{}/Dropbox/REPLICATE-PROJECT/1523841-Quantitative-relationship-between-polarization-differences-a/replication/}
\item \textbf{Status:} tier\_lift\_2026\_04\_26
\end{itemize}

To re-run, change directory into the paper directory above and consult its \texttt{README.md} for the entry-point script. Most replications follow the convention \texttt{replication/run\_*.sh} or \texttt{replication/<subdir>/run.py}.

\end{document}
